\section{Trial setup and the sharp identified interval}

Throughout, generic positive constants may change from line to line, absolute values and Euclidean norms have their usual meanings, and \(O(\cdot)\) and \(o(\cdot)\) denote asymptotic bounds. Arm and label indices range over the sets defined below.

We consider \leanref{S-1}{a randomized causal population} governed by a law \leanref{sym:P}{\ensuremath{P}}. For an overlap margin \leanref{sym:epsilon}{\ensuremath{\varepsilon}}, the score \leanref{sym:E}{\ensuremath{E}} takes values in the interval \leanref{sym:Espace}{\ensuremath{\mathcal E}} given by \([\varepsilon,1-\varepsilon]\). The control and treatment potential outcomes \leanref{sym:Y0}{\ensuremath{Y_0}} and \leanref{sym:Y1}{\ensuremath{Y_1}} take values in the bounded outcome space \leanref{sym:Yspace}{\ensuremath{\mathcal Y=[0,1]}}. The assigned arm is \leanref{sym:A}{\ensuremath{A}}, and the observed outcome is \leanref{sym:Y}{\ensuremath{Y}}. The population score law is \leanref{sym:H}{\ensuremath{H}}; a measurable release rule \leanref{sym:g}{\ensuremath{g}} applied to the score determines the released label \leanref{sym:R}{\ensuremath{R}}; the joint law of the released trial variables is \leanref{sym:Prel}{\ensuremath{P_{\mathrm{rel}}}}. Our target is the average treatment effect \leanref{sym:theta}{\ensuremath{\theta(P)}} for the population law \(P\), defined by \(\theta(P)=\mathbb E_P[Y_1-Y_0]\). For identification, release design, and known-score inference, \(H\) is known to the analyst; in the external-log experiment it is learned from an independent sample of scores.




The population model combines overlap, a specified population score marginal, and score-based randomized assignment. The overlap margin is fixed throughout the analysis.

\begin{assumptionv}{ass:overlap}[Overlap parameter]
The overlap parameter \(\varepsilon\) satisfies \(0<\varepsilon<1/2\).
\end{assumptionv}

The uniform propensity-score overlap condition keeps both assignment probabilities positive throughout the score space \citep{FanShermanShum2014}. It also bounds the reciprocal assignment weights used in the endpoint formulas below.

\begin{assumptionv}{ass:score-marginal}[Score marginal law]
The score \(E\) has law \(H\) under \(P\):
\[
\mathcal L_P(E)=H.
\]
\end{assumptionv}

This specifies the population covariate or propensity marginal used in the data-combination experiment \citep{FanShermanShum2014}. It supplies the score distribution within each released label, even though individual scores are not linked to outcomes in the released data.

\begin{assumptionv}{ass:randomized-assignment}[Score-based treatment assignment]
\[
P(A=1\mid E,Y_0,Y_1)=E
\quad\text{almost surely}.
\]
\end{assumptionv}

Conditional on the score, assignment is Bernoulli and unconfounded with respect to both potential outcomes \citep{FanShermanShum2014}. Thus the score determines how the population score law is weighted in each treatment arm.

The observed outcome and the release are connected to the potential outcomes and score by two further conditions.

\begin{assumptionv}{ass:consistency}[Observed outcome consistency]
\[
Y=AY_1+(1-A)Y_0
\quad\text{almost surely}.
\]
\end{assumptionv}

Potential-outcome consistency identifies the observed outcome with the potential outcome under the assigned arm \citep{Rubin1974}. It connects each observed arm-label outcome distribution to the corresponding potential outcome.

\begin{assumptionv}{ass:deterministic-release}[Deterministic score release]
\[
R=g(E)
\quad\text{almost surely}.
\]
\end{assumptionv}

The deterministic release condition is specific to this analysis: each score is assigned a label by the declared rule. The released label therefore identifies a score cell while retaining uncertainty about the pairing of scores and outcomes within that cell.

\begin{assumptionv}{ass:released-law}[Released data law]
\[
\mathcal L_P(R,A,Y)=P_{\mathrm{rel}}.
\]
\end{assumptionv}

Observed-law compatibility requires a candidate causal law to reproduce the distribution available to the analyst, a standard restriction in data-combination identification \citep{FanShermanShum2014}. Here that distribution includes the label alongside assignment and outcome.

These conditions define the population laws under a release rule and, for a specified released distribution, the compatible laws.

\begin{definitionv}{def:causal-law-class}[Causal-law class \(\mathfrak P(H,g)\)]
The causal-law class \(\mathfrak P(H,g)\) consists of all laws \(P\) satisfying \cref{obj:ass:score-marginal,obj:ass:randomized-assignment,obj:ass:consistency,obj:ass:deterministic-release}.
\end{definitionv}

The class \leanref{sym:Pclass}{\ensuremath{\mathfrak P(H,g)}} in \cref{obj:def:causal-law-class} fixes the score law and release rule while allowing the joint distribution of scores and potential outcomes to vary subject to the stated population conditions.

\begin{definitionv}{def:compatible-law-class}[Compatible-law class \(\mathfrak M(H,g,P_{\mathrm{rel}})\)]
The compatible-law class is
\[
\mathfrak M(H,g,P_{\mathrm{rel}})
=
\bigl\{P\in\mathfrak P(H,g):\mathcal L_P(R,A,Y)=P_{\mathrm{rel}}\bigr\}.
\]
Here \(\mathfrak P(H,g)\) imposes \cref{obj:ass:score-marginal,obj:ass:randomized-assignment,obj:ass:consistency,obj:ass:deterministic-release}, and the released-law condition is \cref{obj:ass:released-law}.
\end{definitionv}

For a fixed released law, \leanref{sym:Mclass}{\ensuremath{\mathfrak M(H,g,P_{\mathrm{rel}})}} in \cref{obj:def:compatible-law-class} is the population class over which the identified set is formed.

We next fix the index sets used to describe released cells. Let \leanref{sym:J}{\ensuremath{J}} denote a positive label count; a later release-design problem uses a positive label budget \leanref{sym:K}{\ensuremath{K}}.

\begin{definitionv}{synth_3}[Released-label alphabets]
For each positive integer \(J\), the released-label alphabet is \(\mathcal R_J=\{1,\ldots,J\}\). In particular, \(\mathcal R_K=\{1,\ldots,K\}\) for a positive label budget \(K\); a release into \(\mathcal R_K\) may use fewer than \(K\) labels.
\end{definitionv}

The alphabet \leanref{sym:Rspace}{\ensuremath{\mathcal R_J}} in \cref{obj:synth_3} includes labels that may have zero probability under a particular release.

\begin{definitionv}{synth_2}[Treatment-arm set]
The treatment-arm set is \(\mathcal A=\{0,1\}\), where \(0\) denotes control and \(1\) denotes treatment.
\end{definitionv}

The arm index \leanref{sym:a}{\ensuremath{a}} ranges over \leanref{sym:Aspace}{\ensuremath{\mathcal A}} as defined in \cref{obj:synth_2}. For each arm-label cell, the observed outcome law and the score law induced by assignment enter the endpoint calculation together.
The endpoint formulas use the generalized quantile \(Q_F(u)=\inf\{x:F((-\infty,x])\ge u\}\) for \(0<u<1\), with support endpoints at \(u=0,1\). This convention also applies when \(F\) has atoms; its formal definition appears in \cref{obj:synth_4}.

\begin{definitionv}{def:arm-cell-primitives}[Arm-cell quantities $p_a,C_r,q_{ar},F_{ar},G_{ar},E_{ar},W_{ar}$]
For each arm \(a\in\mathcal A\) and label \(r\in\mathcal R_J\), define
\[
p_1(e)=e,\qquad p_0(e)=1-e,\qquad
C_r=\{e:g(e)=r\},\qquad
q_{ar}=\int_{C_r}p_a(e)\,H(de).
\]
When \(P_{\mathrm{rel}}(A=a,R=r)>0\), define the observed outcome law
\[
F_{ar}=\mathcal L_{P_{\mathrm{rel}}}(Y\mid A=a,R=r).
\]
For \(q_{ar}>0\), define
\[
G_{ar}(de)=\frac{\mathbf 1\{e\in C_r\}p_a(e)\,H(de)}{q_{ar}},
\qquad E_{ar}\sim G_{ar},
\qquad W_{ar}=\frac{1}{p_a(E_{ar})}.
\]
The endpoint formulas apply to triples \((H,g,P_{\mathrm{rel}})\) satisfying \(P_{\mathrm{rel}}(A=a,R=r)=q_{ar}\) for every arm and label. Their sums include cells with \(q_{ar}>0\), for which \(F_{ar}\) is defined.
\end{definitionv}

For each arm-label cell, the observed outcome law and the score law induced by assignment enter the endpoint calculation together.

\begin{definitionv}{def:mean-endpoints}[Mean endpoints $\underline\mu_a,\overline\mu_a$]
For a triple \((H,g,P_{\mathrm{rel}})\) satisfying \(P_{\mathrm{rel}}(A=a,R=r)=q_{ar}\) for every arm and label, define
\[
\underline\mu_a
=\sum_{r:q_{ar}>0}q_{ar}\int_0^1
Q_{F_{ar}}(u)Q_{W_{ar}}(1-u)\,du,
\qquad
\overline\mu_a
=\sum_{r:q_{ar}>0}q_{ar}\int_0^1
Q_{F_{ar}}(u)Q_{W_{ar}}(u)\,du.
\]
\end{definitionv}

In \cref{obj:def:arm-cell-primitives}, \leanref{sym:p}{\ensuremath{p_a(e)}} is the arm-specific assignment probability, \leanref{sym:C}{\ensuremath{C_r}} is the score cell assigned label \(r\), and \leanref{sym:q}{\ensuremath{q_{ar}}} is its arm-label probability. The observed outcome law \leanref{sym:F}{\ensuremath{F_{ar}}} and the assignment-weighted score law \leanref{sym:G}{\ensuremath{G_{ar}}} are the two marginals available for a positive-mass cell. The reciprocal weight \leanref{sym:W}{\ensuremath{W_{ar}}} converts an arm-specific outcome distribution into a potential-outcome mean; its association with outcomes within the cell determines the endpoint range.

\begin{definitionv}{def:sharp-ate-set}[Sharp ATE interval $I(P_{\mathrm{rel}};H,g)$]
\[
I(P_{\mathrm{rel}};H,g)
=
[\underline\mu_1-\overline\mu_0,\,
 \overline\mu_1-\underline\mu_0].
\]
\end{definitionv}

The generalized quantile \leanref{sym:Q}{\ensuremath{Q_F}} in \cref{obj:synth_4} makes the same endpoint expressions apply to discrete and continuous cell distributions. Pairing outcome and reciprocal-weight quantiles in opposite or matching order gives the lower and upper mean endpoints.

\begin{definitionv}{synth_4}[Generalized quantile]
For a probability law \(F\) with compact support in \(\mathbb R\), define
\[
Q_F(u)=\inf\{x\in\mathbb R:F((-\infty,x])\ge u\},\qquad 0<u<1.
\]
Set \(Q_F(0)=\inf\operatorname{supp}(F)\) and \(Q_F(1)=\sup\operatorname{supp}(F)\). At an atom \(x\), this convention gives \(Q_F(u)=x\) whenever \(F((-\infty,x))<u\le F((-\infty,x])\).
\end{definitionv}

The interval \leanref{sym:Isharp}{\ensuremath{I(P_{\mathrm{rel}};H,g)}} in \cref{obj:def:sharp-ate-set} combines the attainable arm-specific mean endpoints into the fixed-law sharp identified interval for the average treatment effect. The mathematical antecedent is \citet[Theorem 3.2 and eq. (3.1)]{FanShermanShum2014}; the displayed result here is proved under the paper's stated model conditions.

Before reporting that interval for a proposed release, one can check whether the declared score law and released law admit a common causal population. The observable criterion uses the label masses and the treated masses within labels.

\begin{theoremv}{prop:compatibility}[Compatibility of released laws]
Suppose \cref{obj:ass:overlap} holds. Let \(H\) be a probability law on the score interval \(\mathcal E=[\varepsilon,1-\varepsilon]\), let \(g:\mathcal E\to\mathcal R_J\) be a measurable release rule, and let \(P_{\mathrm{rel}}\) be a probability law on \((R,A,Y)\). The compatible causal-law class \(\mathfrak M(H,g,P_{\mathrm{rel}})\) of \cref{obj:def:compatible-law-class} is nonempty if and only if, for every \(r\in\mathcal R_J\),
\[
\begin{aligned}
P_{\mathrm{rel}}(R=r) &= H(C_r),\\
P_{\mathrm{rel}}(A=1,R=r) &= \int_{C_r} e\,H(de).
\end{aligned}
\]
\end{theoremv}

Under \cref{obj:ass:randomized-assignment,obj:ass:deterministic-release}, the score mass of a cell gives its label probability, and score-weighting that mass gives its treated probability. The control probability then follows from the label total. \Cref{obj:prop:compatibility} establishes that these equalities also suffice for a compatible causal law, so an archive can check the coherence of its declared release and randomization law before using the fixed-law interval.
